\section{Second moments of the actual finite financial output}
\label{sec:finite-second-moment}

I strengthen the classical comparator by sampling its fractional financial
output directly. The variance of that output differs from the variance of
the Bernoulli flag in Section~\ref{sec:estimator-coverage}. I do not substitute
the former into the present quantum confidence controller. A quantum mean
estimator that exploits a second-moment promise would need its own complete
construction, confidence proof and cost.

For any real $y\in\R^p$, put
\begin{equation}
 d_i(y)=\mu_i-\overline\mu+(R_i-\overline R)\cdot y,
 \qquad T(y)=\sum_iw_i e^{d_i(y)},\qquad J(y)=\log T(y).
\end{equation}
The weights are positive and sum to one. Hence
$\sum_iw_i d_i(y)=0$ for every $y$, not merely almost surely under a
Gaussian law. Convexity of the exponential gives $T(y)\geq1$. Therefore
$0\leq J(y)\leq T(y)-1$: integrating $1/t\leq1$ from $1$ to $T(y)$
proves the last inequality.

Let $\widetilde Y$ and, for direct evaluation, $\widetilde Z$, have the
independent uniform laws of the actual rounded quantile registers. All
moments in this section first concern these finite laws. They need not
be Gaussian. Let $x$ denote the exact mathematical normalised output of
the chosen financial formula at those register values, before the body
of the financial graph is rounded. For full conditioning,
\begin{equation}\label{eq:second-full}
 0\leq x=P(m_A)-P(m_G)\leq
 \min\{1,L_{h,\lambda}[T(\widetilde Y)-1]\}.
\end{equation}
This follows from $m_A-m_G=J$, the global kernel slope theorem and
$J\leq T-1$. For the random-strike integrand the same proof uses
$L_{\rm CDF}=\varphi(0)/\lambda$, since translation of the CDF argument
does not change its global slope. It holds for every strike-register
value, so averaging over that register preserves the second-moment bound.
Thus
\begin{equation}\label{eq:second-gap}
 \E x^2\leq L^2\{\E T^2-2\E T+1\},
\end{equation}
with $L=L_{h,\lambda}$ for full conditioning and $L=L_{\rm CDF}$ for
random-strike conditioning. It is essential to retain the negative
$-2\E T$ term with a lower, rather than an upper, enclosure of $\E T$.

I evaluate these moments without a joint payoff table. Define the scalar
moment-generating function $M_{b,q}(t)=\E e^{t\widetilde Y_j}$ for the
common marginal with latent width $b$ and fractional width $q$. Independence
and multiplication of finite sums give, for any real $c$ and vector $v$,
\begin{equation}\label{eq:finite-exponential}
 \E e^{c+v\cdot\widetilde Y}
 =e^c\prod_{j=1}^pM_{b,q}(v_j).
\end{equation}
Writing $s_i=R_i-\overline R$ and $a_i=\mu_i-\overline\mu$ therefore yields
\begin{align}\label{eq:finite-t-moments}
 \E T&=\sum_iw_i e^{a_i}\prod_jM_{b,q}(s_{ij}),\\
 \E T^2&=\sum_{i,k}w_iw_k e^{a_i+a_k}
                  \prod_jM_{b,q}(s_{ij}+s_{kj}).
\end{align}
These identities are exact for the stipulated finite law. Replacing
$M_{b,q}(t)$ by $e^{t^2/2}$ gives the separate continuous Gaussian
comparison. That substitution is not made in the finite certificate.

An additional bound follows from the positive arithmetic--geometric gap.
Set $B=\sum_iw_i e^{\mu_i+R_i\cdot\widetilde Y}$ and
$C=e^{\overline\mu+\overline R\cdot\widetilde Y}$, so $B\geq C$.
The unit Lipschitz payoff gives $F_H(A)-F_H(G)\leq A-G$.
For direct evaluation this proves
\begin{equation}\label{eq:second-direct}
 \E x^2\leq \frac{K^2}{H^2}M_{b_Z,q}(2\lambda)
                       \E(B-C)^2.
\end{equation}
For full conditioning, taking the continuous inner Gaussian expectation
first gives $Hx\leq K e^{\lambda^2/2}(B-C)$, and hence
\begin{equation}\label{eq:second-condition-AG}
 \E x^2\leq\frac{K^2e^{\lambda^2}}{H^2}\E(B-C)^2.
\end{equation}
The two inner factors are different: the square of
$\E e^{\lambda Z}$ equals $e^{\lambda^2}$, whereas
$\E e^{2\lambda Z}=e^{2\lambda^2}$. For a finite direct register the
last factor is its actual $M_{b_Z,q}(2\lambda)$.
Expanding the square gives
\begin{align}
 \E(B-C)^2={}&
 \sum_{i,k}w_iw_k e^{\mu_i+\mu_k}
      \prod_jM_{b,q}(R_{ij}+R_{kj})\\
 &-2\sum_iw_i e^{\mu_i+\overline\mu}
      \prod_jM_{b,q}(R_{ij}+\overline R_j)
 +e^{2\overline\mu}\prod_jM_{b,q}(2\overline R_j).
\end{align}
The cross term again requires a lower enclosure. I use the minimum of the
applicable proved bounds and $\E x^2\leq1$.

\subsection{A scalar certificate for the dyadic marginal}

The installed quantile constructor rounds the same floating seed
$s_k=\operatorname{NormalDist.inv\_cdf}((k+1/2)/2^b)+5$ to each chosen
fractional width and imposes exact mirror symmetry. This description
defines its actual integer law; it does not assert that the floating seed
is the exact Gaussian quantile. I retain a reference table at width
$q_*=48$. The largest scaled seed is below $2^{53}$, and scaling a finite
binary floating number by the powers of two used here is exact. Nearest
integer rounding consequently gives the pointwise displacement
\begin{equation}\label{eq:rounding-marginal}
 |\widetilde y_{k,q}-\widetilde y_{k,q_*}|
 \leq \delta_q:=\tfrac12(2^{-q}+2^{-q_*}).
\end{equation}
The reference and production mirror operations preserve this bound on the
negative half. Reproduction compares their integer maps with the installed
constructor, rather than relying only on a floating-distribution name.

Symmetry gives $M_{b,q_*}(t)=\E\cosh(t\widetilde Y_*)$. Exact integer
power sums supply every even moment of this dyadic table. For $m=64$,
\begin{equation}
 P_m(t)=\sum_{k=0}^{m}\frac{t^{2k}\E\widetilde Y_*^{2k}}{(2k)!}
 \leq M_{b,q_*}(t).
\end{equation}
If $a=|t|\max_k|\widetilde y_{k,q_*}|$, the first omitted term is at
most $a^{2m+2}/(2m+2)!$. Every subsequent term ratio is at most
$\rho=a^2/[(2m+3)(2m+4)]<1$. Summing a geometric majorant proves
\begin{equation}
 0\leq M_{b,q_*}(t)-P_m(t)
 \leq\frac{a^{2m+2}}{(2m+2)!(1-\rho)}.
\end{equation}
The retained coefficient and support checks make this last bound smaller
than $10^{-120}$. Since
$e^{-|t|\delta_q}e^{t\widetilde Y_*}\leq
 e^{t\widetilde Y_q}\leq
 e^{|t|\delta_q}e^{t\widetilde Y_*}$ pointwise, taking expectations
encloses the production MGF by multiplying the reference lower and upper
bounds by the respective exponential factors. All power sums and
coefficient inputs are exact integers or rational numbers. The Taylor sum
uses correctly rounded decimal operations at 220 significant digits and an
aggregate $10^{-100}$ outward envelope. Its at most 65 moment terms have
magnitudes below $5^{128}<10^{90}$; the coefficient magnitude is at most one,
and the remaining moment, density, exponential and budget intermediates
have magnitude below $10^{100}$. Fewer than $10^6$ correctly rounded
operations contribute to any retained scalar bound, so their conservative
accumulated absolute roundoff is below $10^{-113}$, before the larger
$10^{-100}$ envelope. I also apply that envelope at subsequent product and
moment boundaries. The finite set of operations, their magnitudes and the
Taylor remainder premises are checked independently. No mean-error certificate is treated
as a second-moment or pointwise distribution certificate.

\subsection{Transfer to the implemented output and confidence budget}

Let $X$ be the actual implemented fractional output. The graph recurrence
has a pointwise currency error bound $a_{\rm body}$ that includes a
conservative allowance for quantile-register error. I can set the ideal
quantile node equal to its actual rounded value. Its initial error then
equals zero, which is no larger than the retained allowance, and its value
lies in the certified quantile range. Every subsequent propagation rule is
monotone in the incoming nonnegative error allowances. The original domain
and overflow enclosures continue to contain the resulting subsets.
Consequently the same retained bound proves
\begin{equation}
 |X-x|\leq a_{\rm body}/H
\end{equation}
at every finite address. Generator enclosures, interpolation, constant
rounding, clamps and product truncation remain included in that bound.

If $M_2$ is a proved bound for $\E x^2$, Cauchy--Schwarz gives
$\E[x|X-x|]\leq\sqrt{M_2}\,a_{\rm body}/H$. Expanding
$X^2\leq(x+|X-x|)^2$ therefore proves
\begin{equation}\label{eq:finite-second-transfer}
 \E X^2\leq\min\{1,(\sqrt{M_2}+a_{\rm body}/H)^2\}=:M_{2,+}.
\end{equation}
This is the $L^2$ triangle inequality written out for the present finite
law. The moment transfer does not use the distribution bias budget. That
budget still enters the final target-price guarantee through the unchanged
deterministic mean-bias certificate.

Since $0\leq X\leq1$, its variance is at most $1/4$. Also
$\operatorname{Var}(X)=\E X^2-(\E X)^2\leq M_{2,+}$.
Combining this with the already proved mean-prior variance bound
$v_{\rm mean}$ gives
\begin{equation}
 v_+=\min\{1/4,M_{2,+},v_{\rm mean}\}.
\end{equation}
For independent exact draws of $X$, the previously derived Bernstein
inequality consequently gives the sufficient fixed sample count
\begin{equation}\label{eq:finite-variance-samples}
 N_B=\left\lceil
 \frac{(2v_++2t/3)\log(2/\alpha)}{t^2}\right\rceil,
 \qquad t=0.005/H.
\end{equation}
I also retain the smaller applicable Hoeffding count. The count is fixed
before sampling; choosing the certified variance or prefix through
deterministic model calculations does not invalidate coverage. At
$\alpha=0.01$, the finite mean has error at most $0.005$ in undiscounted
currency with probability at least $0.99$. The retained deterministic bias
of at most $0.005$ and $e^{-r}\leq1$ give absolute price error at most
$0.01$. For one preselected price per case, allocating
$\alpha=0.01/27$ and applying the union bound gives simultaneous
$99$ per cent coverage over the 27 prices. This does not certify every
retained candidate output simultaneously.

The financial inequalities, exponential moment factorisation, exact dyadic
sums, Taylor majorant, $L^2$ transfer and concentration inequality are
established mathematics. Their combined parameter certificate is a
project-specific argument. I have not established theoretical originality
or a new quantum primitive from this calculation.
